% CODEX REPORT::LEAD::FINAL-LINK2-ARC::2026-10-08
\documentclass[11pt]{article}
\usepackage[T1]{fontenc}
\usepackage[margin=0.80in]{geometry}
\usepackage{amsmath,amssymb,amsthm,booktabs,array,mathtools,xcolor,xurl,hyperref}
\hypersetup{colorlinks=true,linkcolor=blue!55!black,citecolor=blue!55!black,urlcolor=blue!55!black}
\setlength{\emergencystretch}{3em}
\setlength{\parskip}{0.35em}
\setlength{\parindent}{0pt}
\newtheorem{theorem}{Theorem}[section]
\newtheorem{lemma}[theorem]{Lemma}
\newtheorem{corollary}[theorem]{Corollary}
\newtheorem{definition}[theorem]{Definition}
\newtheorem{proposition}[theorem]{Proposition}
\newcommand{\eps}{\varepsilon}
\newcommand{\SAT}{\operatorname{SAT}}
\newcommand{\VALUE}{\mathsf{VALUE}}
\newcommand{\REC}{\mathsf{REC}}
\newcommand{\rootT}{i_0}
\title{FINAL-LINK2-ARC:\\Typed Double Wrappers, Potential-Collapsed Epochs, and Mixed Exit Channels}
\author{Lead ARC audit}
\date{8 October 2026}

\begin{document}
\maketitle

\begin{abstract}
This report isolates a genuine relaxation of the missing universal
source-only cover.  A valid final compiler need not recursively cover every
input at every raw step.  It may first execute a charged, source-checkable
administrative epoch, terminate some branches by directly obtained
\(\VALUE\) answers, and send the remaining branches through typed recursive
channels.  Same-scale raw steps may occur inside the epoch when a proved
potential bounds their total number, depth, acquisition cost, replay cost,
and live state.  After each epoch is collapsed, however, every feasible
infinite return cycle must still contract.  The exact typed time budgets are
\[
 C_i\ge R_i,\qquad C_i\ge w_e+\theta_e C_j,
\]
and the exact streamed-space budgets are
\[
 \rho_i\ge P_i,\qquad \rho_i\ge h_e,
 \qquad \rho_i\ge\theta_e\rho_j.
\]
For a two-wrapper cycle \(A\to B\to A\), closed forms are derived below,
including the case of a same-scale outer phase.  A conditional charged point
\((C_A,\rho_A)=(6/5,1/2)\) lies below the Williams line, but no recovered
construction realizes it for every general CNF.  An exact potential theorem
also yields a conditional point \((1,3/4)\).  A 9,598-check rational audit
verifies the formulas and rejects unit-scale cycles, omitted fanout, omitted
stack depth, parent-scale \(\VALUE\) queries, and mismatched typewise
time-space witnesses.  These results sharpen the final target.  They do not
complete it, and \(P\) versus \(NP\) remains \textsc{Open}.
\end{abstract}

\paragraph{Run identity and immutable scope.}
The run identifier is \texttt{FINAL-LINK2-ARC}.  This report is an immutable
successor to the FINAL-LINK1 analysis in \texttt{main8.tex}; it does not edit
\texttt{main4.tex}, \texttt{main5.tex}, or \texttt{main8.tex}.  The words
\textsc{Proved/Exact}, \textsc{Conditional}, \textsc{Computational}, and
\textsc{Open} are used literally.  An exponent identity proved from stated
recurrences is \textsc{Proved/Exact}.  An algorithmic consequence depending
on a not-yet-constructed all-input SAT mechanism is \textsc{Conditional}.
Finite experiments are \textsc{Computational}.  The constructional gap and
the separation claim are \textsc{Open}.

\section{Outcome: what can actually be relaxed}

Set
\[
 \gamma=2\cos(\pi/7)=1.8019377358\ldots.
\]
The endpoint used by the proof spine forbids one deterministic, uniform,
exact random-access SAT machine with simultaneous exponents \((C,\rho)\)
satisfying \(C+\rho<\gamma\) \cite{Williams2010,FortnowEtAl2005}.

The all-input quantifier cannot be removed from a worst-case theorem.  It can
be moved.  Instead of requiring every input to have a source-only recursive
cover at every raw wrapper step, it is enough that every input receive one
charged exact \emph{epoch plan} whose terminal channels are of the following
two kinds:
\begin{align*}
 \VALUE&:\ \text{a correct Boolean value obtained within the local budget},\\
 \REC&:\ \text{a typed recursive exit of certified scale.}
\end{align*}
The epoch may contain same-scale administrative steps and may choose different
mechanisms for different source-checkable input types.  Its complete cost is
collapsed into local floors and recursive edge weights.  The collapsed type
graph must then have favorable Bellman budgets, and every feasible recurrent
cycle must shrink.  Thus the relaxation is logical and useful, but it does
not turn worst-case coverage into average coverage.

\begin{theorem}[Mixed-epoch final-link target, \textsc{Conditional}]
Suppose a fixed uniform compiler gives every CNF input a source-checkable,
exact, finite epoch plan.  Every terminal leaf is either a charged \(\VALUE\)
leaf or a charged typed \(\REC\) exit.  Suppose all discovery, failed
portfolio members, fanout, generation, acquisition, source rereads, virtual
access, replay, combination, retained state, and live workspace satisfy the
typed ledgers in Sections~\ref{sec:ledger} and~\ref{sec:potential}.  If every
feasible reachable recurrent cycle contracts and the least root budgets obey
\[
 C_{\rootT}^*+\rho_{\rootT}^*<\gamma,
\]
then the compiler supplies the missing finite entry into the forbidden
region.  Combined with the established endpoint and a finite polynomial SAT
seed under \(P=NP\), it yields a contradiction.
\end{theorem}

The deduction from the premises is \textsc{Proved/Exact}.  The existence of
such a compiler for all CNFs is \textsc{Open}.  This is the decisive status
split.

\section{Fully charged typed ledger}\label{sec:ledger}

Let \(I\) be a fixed finite set of wrapper types.  An edge
\(e:i\to j\) represents one class of recursive children.  Its scale is
\(0\le\theta_e\le1\).  The following separation of charges prevents a
common ambiguity.

\begin{description}
\item[Recursive multiplier \(w_e\).]
If there are \(N^{f_e+o(1)}\) child calls and simulating one unit of child work
requires \(N^{u_e+o(1)}\) interface, translation, or repeated-access work,
then \(w_e=f_e+u_e\).  The recursive term is
\(N^{w_e+\theta_e C_j+o(1)}\).

\item[Local time floor \(R_i\).]
This is the maximum exponent of setup, type classification, discovery, all
failed portfolio trials, batch acquisition, one-time generation, source
rereads, materialization, exact verification, \(\VALUE\)-leaf work,
combination, output, and replay that is not multiplicative per unit of child
solver work.  For example, \(N^{f_e}\) children with one-time generation
cost \(N^{g_e}\) contribute \(f_e+g_e\) to \(R_i\), not automatically to
\(w_e\).  If replay occurs on every simulated child step, it belongs in
\(u_e\); if it occurs once per child, it belongs in \(R_i\).

\item[Local live-space floor \(P_i\).]
This includes the source representation or handle, type-classifier state,
generator state, administrative stack, virtual decoder, batch tables,
materialized plan records, explicit child buffers, and \(\VALUE\) workspaces.

\item[Retained edge state \(h_e\).]
This is the parent state that remains live while the recursive child runs.
Sequential evaluation is assumed.  If children run in parallel, their
aggregate memory, including fanout, must replace this streamed bound.
\end{description}

The time and space Bellman inequalities are
\begin{align}
 C_i&\ge R_i, & C_i&\ge w_e+\theta_eC_j
     &&(e:i\to j),\label{eq:timebellman}\\
 \rho_i&\ge P_i, & \rho_i&\ge h_e,&
 \rho_i&\ge\theta_e\rho_j
     &&(e:i\to j).\label{eq:spacebellman}
\end{align}
For ordinary explicit children, \(P_i\) must contain their live encoded size,
normally at least \(\theta_e\).  A virtual child may omit that buffer only
when a bit-level decoder is specified.  Decoder state belongs in \(P_i\),
one-time decoder work belongs in \(R_i\), and per-query translation belongs
in \(w_e\).  Calling an assumed SAT solver merely to decide which edge or
\(\VALUE\) leaf to use produces a solver term at that call's actual scale.
In particular, a parent-scale answer has \(\theta=1\); it is not local free
work.

\begin{theorem}[Bellman path characterization, \textsc{Proved/Exact}]
For the declared exponent recurrence, the least finite time budget is
\[
 C_i^*=\sup_P\left(W_P+\Theta_P R_{\operatorname{end}(P)}\right),
\]
where the supremum is over finite directed paths beginning at \(i\),
\[
 \Theta_P=\prod_{e\in P}\theta_e,
 \qquad
 W_P=\sum_{k=0}^{|P|-1}
       \left(\prod_{\ell<k}\theta_{e_\ell}\right)w_{e_k}.
\]
The least streamed-space budget is the analogous maximum of every local or
retained-state floor transported back by the prefix products.  Equations
\eqref{eq:timebellman}--\eqref{eq:spacebellman} are therefore necessary and
sufficient within this recurrence abstraction.  They are not lower bounds on
all possible SAT algorithms.
\end{theorem}

\begin{proof}
Repeated substitution of the Bellman inequalities gives the lower bound for
every finite path.  Conversely, the path supremum dominates every local floor
and, after prefixing an edge, dominates \(w_e+\theta_eC_j^*\).  It is thus a
pre-fixed point and is no larger than any other pre-fixed point.  The space
proof is identical with zero additive edge weights and with each encountered
local or retained floor transported by the preceding scale product.
\end{proof}

A finite portfolio of exact mechanisms can be represented by additional edge
classes and local terms.  Finite sums become maxima of polynomial exponents.
The maximum must include every failed mechanism tested before success, not
only the successful member.

\section{Complete return cycles}

For a directed cycle \(Z=e_0e_1\cdots e_{m-1}\) starting at \(i_0\), define
\[
 \Theta_Z=\prod_{k=0}^{m-1}\theta_{e_k},
 \qquad
 W_Z=\sum_{k=0}^{m-1}
 \left(\prod_{\ell<k}\theta_{e_\ell}\right)w_{e_k}.
\]
Substitution around the cycle gives the exact necessary inequality
\[
 C_{i_0}\ge W_Z+\Theta_Z C_{i_0}.
\]
Hence, if \(\Theta_Z<1\),
\begin{equation}
 C_{i_0}\ge \frac{W_Z}{1-\Theta_Z}.
 \label{eq:cyclethreshold}
\end{equation}
If \(\Theta_Z=1\) and \(W_Z>0\), no finite time budget exists.  If
\(\Theta_Z=1\) and \(W_Z=0\), finite fixed points may exist, but a large
inherited solver exponent circulates unchanged.  Such a cycle cannot by
itself establish finite entry from an unknown polynomial seed.

\begin{theorem}[Cycle contraction criterion, \textsc{Proved/Exact}]
Assume nonnegative weights and scales in \([0,1]\).  A finite typed return
system suppresses every finite inherited exponent vector if every feasible
reachable directed cycle has \(\Theta_Z<1\).  Conversely, a feasible
reachable cycle with \(\Theta_Z=1\) prevents generic suppression of an
arbitrarily large inherited exponent.  Cycle thresholds alone are not
sufficient for the endpoint: all local floors and paths into them must also
be included through the least Bellman solution.
\end{theorem}

\begin{proof}
In a finite graph, sufficiently long paths decompose into a bounded simple
prefix plus cycles.  Strict contraction of every cycle makes the scale
product along every infinite path tend to zero, so the inherited boundary
term disappears and the path supremum converges to the least Bellman budget.
If a cycle has product one, every edge on it has scale one.  Starting with a
constant vector \(K\) and following that cycle leaves the boundary term at
least \(K\); positive cycle weight additionally makes it grow.
\end{proof}

Only feasible cycles count.  If a source-checkable type invariant proves two
edges mutually incompatible on one execution, the type automaton should
encode that invariant before cycles are enumerated.  Merely hoping that a
syntactic cycle is not realized is not a proof.

\section{Exact typed double-wrapper region}

Consider two types.  The edge \(A\to B\) has scale \(\alpha\) and charged
recursive multiplier \(x\); the edge \(B\to A\) has scale \(\beta\) and
multiplier \(y\).  Let \(R_A,R_B\) be the complete time floors.  Define
\[
 H_A=\max\{P_A,h_{AB}\},\qquad
 H_B=\max\{P_B,h_{BA}\}.
\]

\begin{theorem}[Two-wrapper closed form, \textsc{Proved/Exact}]
If \(\alpha\beta<1\), the least budgets are
\begin{align}
 C_A^*&=\max\left\{R_A,\ x+\alpha R_B,
             \frac{x+\alpha y}{1-\alpha\beta}\right\},\label{eq:CA}\\
 C_B^*&=\max\left\{R_B,\ y+\beta R_A,
             \frac{y+\beta x}{1-\alpha\beta}\right\},\label{eq:CB}\\
 \rho_A^*&=\max\{H_A,\alpha H_B\},\qquad
 \rho_B^*=\max\{H_B,\beta H_A\}.\label{eq:rhoAB}
\end{align}
For root type \(A\), the exact acceptance inequality for this recurrence is
\begin{equation}
 \boxed{
 \max\left\{R_A,x+\alpha R_B,
       \frac{x+\alpha y}{1-\alpha\beta}\right\}
 +\max\{H_A,\alpha H_B\}<\gamma.}
 \label{eq:doubleaccept}
\end{equation}
\end{theorem}

\begin{proof}
Substituting
\(C_B=\max\{R_B,y+\beta C_A\}\) into the equation for \(C_A\) yields
\[
 C_A=\max\{R_A,x+\alpha R_B,x+\alpha y+alpha\beta C_A\}.
\]
The least solution is \eqref{eq:CA}; symmetry gives \eqref{eq:CB}.  Applying
the same substitution to the zero-additive space equations gives
\eqref{eq:rhoAB}.
\end{proof}

The particularly relevant same-scale outer wrapper has \(\alpha=1\) and
\(\beta=b<1\).  Its root criterion is
\begin{equation}
 \boxed{
 \max\left\{R_A,x+R_B,\frac{x+y}{1-b}\right\}
 +\max\{H_A,H_B\}<\gamma.}
 \label{eq:samescaleouter}
\end{equation}
This shows exactly how a double wrapper may use one noncontracting phase.  It
does not permit an unresolved current-scale cycle.

As a rational positive control, take
\[
 \alpha=1,\quad\beta=\tfrac12,\quad
 x=\tfrac15,\quad y=\tfrac25,\quad
 R_A=R_B=1,\quad H_A=H_B=\tfrac12.
\]
Then
\[
 C_A^*=\tfrac65,\quad C_B^*=1,
 \quad\rho_A^*=\rho_B^*=\tfrac12,
 \quad C_A^*+\rho_A^*=\tfrac{17}{10}<\gamma.
\]
The arithmetic is \textsc{Proved/Exact}; its realization by a general-CNF
wrapper is \textsc{Conditional} and presently \textsc{Open}.

\section{Potential-collapsed administrative epochs}\label{sec:potential}

This is the strongest provable relaxation found in the present run.  It
allows same-scale raw work without allowing an uncharged same-scale solver
cycle.

\begin{theorem}[Mass-potential epoch bound, \textsc{Proved/Exact}]
Fix an input of length \(N\).  Suppose a compiler constructs a finite rooted
administrative plan tree.  Each node \(v\) has a positive, source-checkable
mass \(\mu(v)\).  Assume
\begin{align*}
 \mu(\mathrm{root})&\le N^{\kappa+o(1)},\\
 \sum_{u\text{ child of }v}\mu(u)&\le\mu(v),\\
 \mu(v)&\ge N^{-\xi+o(1)}\quad\text{for every internal node},\\
 D(N)&\le N^{d+o(1)}\quad\text{for the tree depth}.
\end{align*}
If internal-node processing costs \(N^{u+o(1)}\), its total time exponent is
at most
\begin{equation}
 r_{\mathrm{internal}}=\kappa+\xi+d+u.\label{eq:internalpotential}
\end{equation}
If every terminal exit in recursive class \(e\) has mass at least
\(N^{-\xi_e+o(1)}\), then the number of such exits has exponent at most
\begin{equation}
 f_e=\kappa+\xi_e.\label{eq:exitpotential}
\end{equation}
If each exit has one-time acquisition or generation cost exponent \(g_e\),
then \(R_i\) contains \(f_e+g_e\).  If every unit of recursive solver work
incurs access/replay exponent \(u_e\), then the Bellman edge weight is
\begin{equation}
 w_e=f_e+u_e=\kappa+\xi_e+u_e.\label{eq:edgepotential}
\end{equation}
A \(\VALUE\) leaf class \(q\) with minimum leaf mass
\(N^{-\xi_q+o(1)}\) and per-leaf decision cost \(N^{v_q+o(1)}\)
contributes \(\kappa+\xi_q+v_q\) to \(R_i\).
\end{theorem}

\begin{proof}
At any fixed depth, conservation gives total mass at most the root mass.
Therefore the number of internal nodes at that depth is at most
\(N^{\kappa+\xi+o(1)}\).  Multiplying by depth and per-node work proves
\eqref{eq:internalpotential}.  All leaves form an antichain.  Repeated use of
the conservation inequality gives total leaf mass at most the root mass, so
class \(e\) has at most \(N^{\kappa+\xi_e+o(1)}\) exits, with no extra depth
factor.  The remaining assertions follow by multiplying the number of exits
by the declared per-exit or per-unit costs.
\end{proof}

The potential certificate is not free.  Computing its types, weights, or
local inequalities belongs in \(R_i\) and \(P_i\).  A certificate whose
verification enumerates the whole plan is charged at that enumeration cost.
For streamed depth-first evaluation, if every live frame has exponent \(h\),
the administrative stack has exponent \(h+d\).  A safe local space floor is
therefore
\begin{equation}
 P_i\ge\max\{\sigma_i,h+d,max_e s_e,
     \max_{e\text{ explicit}}\theta_e,p_{\mathrm{plan}}\},
 \label{eq:potentialspace}
\end{equation}
where \(\sigma_i\) is source/summary state, \(s_e\) is generator or decoder
state, and \(p_{\mathrm{plan}}\) is zero only if the plan is truly streamed.

Same-scale steps may appear in this administrative tree because the potential
and rank force the epoch to end.  After the epoch is collapsed, complete
macro return cycles must still satisfy the contraction criterion.  A raw
unit-scale loop with no well-founded rank cannot be relabeled as amortized.

One conditional arithmetic point is
\[
 \kappa=\tfrac16,\quad\xi=\xi_e=\tfrac1{12},
 \quad d=0,\quad u=\tfrac34,\quad u_e=0,
 \quad\theta_e=\tfrac34.
\]
Equations~\eqref{eq:internalpotential}--\eqref{eq:edgepotential} give
\(R=1\), \(w=1/4\), and \(C^*=1\).  Explicit child storage gives
\(P\ge3/4\), so \((C^*,P)=(1,3/4)\) has sum \(7/4<\gamma\).  This is a
\textsc{Conditional} target.  No all-CNF mass certificate with these
parameters was recovered.

\section{Typed nonuniformity and the same-realization guard}

A finite classifier may choose distinct mechanisms for distinct types while
the overall algorithm remains uniform.  The classifier must be total,
computable without the desired SAT answer, and fully charged.  If \(I_0\) is
the set of possible root types, the uniform worst-case algorithm has
\begin{equation}
 C_{\mathrm{global}}=\max_{i\in I_0}C_i^*,\qquad
 \rho_{\mathrm{global}}=\max_{i\in I_0}\rho_i^*.
 \label{eq:globalmaxima}
\end{equation}
The endpoint requires
\begin{equation}
 \boxed{\max_i C_i^*+\max_i\rho_i^*<\gamma.}
 \label{eq:samerealization}
\end{equation}
It is not enough that \(C_i^*+\rho_i^*<\gamma\) separately for every type.
For example, the two pairs
\[
 (17/10,1/20),\qquad(4/5,1)
\]
each have coordinate sum below \(\gamma\), but their global maxima have sum
\(17/10+1=27/10\).  The worst time and worst space may occur on different
inputs; complexity functions take the maxima independently.

A fixed finite portfolio can be tried sequentially, but all failed trials are
part of \(R_i\).  If every member is exact and total, selecting one based on
a cheap structural certificate is permitted.  Selecting the favorable member
using a parent SAT answer is circular.  Parallel execution does not change a
time exponent for a fixed finite portfolio, but it may increase live space;
for a growing portfolio both its count and simultaneous state must be charged.

\section{Type potentials redistribute cost but create no slack}

Let \(\pi_i\) be a credit attached to type \(i\), and write
\(D_i=C_i-\pi_i\).  The reweighted edge is
\[
 w'_e=w_e+\theta_e\pi_j-\pi_i
 \qquad(e:i\to j).
\]

\begin{proposition}[Cycle invariance, \textsc{Proved/Exact}]
For a cycle \(Z\) based at \(i_0\),
\[
 W'_Z=W_Z-(1-\Theta_Z)\pi_{i_0},
\]
and therefore
\[
 \frac{W'_Z}{1-\Theta_Z}+\pi_{i_0}
 =\frac{W_Z}{1-\Theta_Z}.
\]
\end{proposition}

\begin{proof}
Substitute the definition of \(w'_e\) into the discounted cycle sum.  Every
intermediate potential cancels; the initial potential occurs with
coefficients \(-1\) and \(\Theta_Z\).
\end{proof}

Type potentials can expose an amortized proof and move apparent cost between
phases.  They cannot manufacture endpoint margin.  Any claimed improvement
that appears only after reweighting but disappears when the root credit is
restored fails this invariant.

\section{Mixed \texorpdfstring{\(\VALUE\)}{VALUE} and recursive channels}

Direct \(\VALUE\) terminals are a real relaxation.  An input type whose
answer is obtained by a charged local rule needs no recursive cover.  More
generally, a Boolean AND/OR plan may mix \(\VALUE\) leaves and recursive
leaves.  Exactness fixes the worst cases: an OR plan may need all children on
an UNSAT input, and an AND plan may need all children on a SAT input.
Short-circuiting helps particular inputs but cannot erase worst-case fanout.

A \(\VALUE\) leaf produced by running the supplied solver on the parent is
not a terminal constant.  It is a solver edge with \(\theta=1\).  Likewise,
writing ``answer later'' on a virtual child does not charge the later source
queries.  The mixed-channel theorem helps precisely when the \(\VALUE\)
acquisition or structural recognition is already within \(R_i\).

The final all-input requirement can now be stated narrowly: every input must
receive a favorable \emph{terminal-or-recursive epoch}.  It need not receive
a purely recursive cover.  Average success over an input distribution,
average success over sampled hashes, and amortization across independent user
inputs do not satisfy the deterministic worst-case endpoint.  Randomized
expected bounds also require a model-matched randomized endpoint or a proved
derandomization before they can enter this spine.

\section{Recovered project work and what it proves}

The local corpus was searched for ARC, ratchet, final-crossing, return-cycle,
multichannel, stream, batch, and resource-bound statements.  The relevant
pieces reconcile as follows.

\begin{center}
\small
\begin{tabular}{p{0.22\linewidth}p{0.18\linewidth}p{0.50\linewidth}}
\toprule
Recovered component & Status here & Effect on FINAL-LINK2 \\
\midrule
MULTICHANNEL1 & \textsc{Proved/Exact} accounting & Gives the finite-channel
max-affine fixed point when every channel shrinks. \\
STREAMCOMP1 & \textsc{Conditional} interface ledger & Charges virtual access,
generation, and typed recursive space, provided a concrete decoder realizes
the declared interface. \\
ORBITBATCH1 & \textsc{Conditional} batch ledger & Separates shared setup and
classification from recursive work; does not prove a general-SAT batch. \\
CYCLEGAIN1 and finite return cycles & \textsc{Proved/Exact} arithmetic,
\textsc{Open} SAT successor & Permit individually expanding or same-scale
phases when the complete feasible return cycle contracts. \\
MACRO-LINEAGE-SPACE1 & \textsc{Proved/Exact} in scope & Bounds route-index bits;
does not bound formulas, caches, decoders, or stored plans. \\
Precommit and ratchet records & \textsc{Proved/Exact} binding in scope & Provide
identity, reseparation, or reacquisition obligations, not a resource-decreasing
potential for unrestricted CNF. \\
Forward double wrapper & \textsc{Proved/Exact} semantics,
\textsc{Open} favorable resources & Supplies exact AND/OR plans and small
leaves, but the current implementations store formulas, plan DAGs, and cache
keys and do not meet \eqref{eq:samescaleouter}. \\
Cofinal ARC seed hitting & \textsc{Proved/Exact} implication,
\textsc{Open} SAT hit & Shows a universal common contraction is not logically
necessary; it still needs one charged SAT-specific finite route. \\
\bottomrule
\end{tabular}
\end{center}

No recovered \textsc{Proved/Exact} item supplies an all-input potential
certificate, mixed epoch, or typed double-wrapper realization whose complete
root budgets satisfy the endpoint.  In particular, semantic ratcheting is not
resource contraction: a terminal jump can postpone all hard work until the
last step.

\section{Primary-literature clues}

Measure-and-conquer analyses support the potential-epoch form: nonuniform
branch progress can be amortized with a carefully chosen measure, but only
after each branch rule and its work are proved against that measure
\cite{FominGrandoniKratsch2009}.  This motivates Section~\ref{sec:potential};
it does not provide the missing unrestricted-CNF measure.

Recursive SAT backdoors and backdoor depth show that different variables may
be used in different recursively created components and can yield tractable
structured formula families \cite{MaehlmannSiebertzVigny2021,
DreierOrdyniakSzeider2022}.  Backdoor DNFs provide exact disjunctive covers
whose terms move a formula into a tractable base class
\cite{OrdyniakSchidlerSzeider2021}.  These are strong clues for typed and mixed
plans, but their positive guarantees are parameterized or structural.  They
do not prove that every CNF has endpoint-compatible depth, fanout,
recognition, access, and live space.

Williams's self-improvement work is a precedent for finite bootstrapping from
a weak algorithmic hypothesis, but it does not supply the SAT wrapper needed
here \cite{Williams2024}.  The literature therefore supports the accounting
architecture and several restricted mechanisms, not the open all-input link.

\section{Computational audit}

The executable audit is
\path{pnp_spines/final_link2_typed_audit.py}.  It installs a
25-second alarm and completed in 1.57 seconds.  It uses exact rational
arithmetic for all ledger identities; floating point is used only to compare
rational totals with \(\gamma\).  Its JSON record is
\path{pnp_spines/final_link2_typed_audit_results.json}.

The run performed 9,598 checks.  It compared the two-wrapper closed form with
exact policy enumeration on 9,375 small rational parameter choices and
tested 200 random three-type contracting cycles against long Bellman
iteration.  The candidate \((6/5,1/2)\) passed the numerical endpoint.

The hostile controls all fired as intended:
\begin{description}
\item[Positive unit-product cycle.]
Exact policy enumeration found no finite Bellman budget.
\item[Zero-weight unit-product cycle.]
An inherited exponent \(K=7\) remained \(7\) after 100 returns.
\item[Same-realization mismatch.]
Two individually favorable typed pairs produced the global sum \(27/10\).
\item[Omitted branching.]
A binary same-scale tree of depth \((1/4)\log_2N\) produced exit exponent
\(1/4\), rejecting a zero-fanout ledger.
\item[Omitted live stack.]
Frame exponent \(3/10\) and depth exponent \(1/5\) produced space exponent
\(1/2\), not \(3/10\).
\item[Parent \(\VALUE\) query.]
A scale-one zero-weight query preserved the inherited exponent.
\item[Potential reweighting.]
The discounted cycle threshold was invariant after restoring the root credit.
\end{description}

These are \textsc{Computational} checks of the algebra and the audit program.
They are not empirical evidence that the missing SAT construction exists.

\section{Proposed manuscript changes}

No proof-spine manuscript was edited in this run.  The following changes are
recommended for a reconciled successor to \texttt{main8.tex}.

\begin{description}
\item[Replace source-only coverage by mixed-epoch coverage.]
State that each input may terminate through a charged \(\VALUE\) rule or
leave through a typed recursive channel.  Preserve the all-input quantifier
at the epoch level.

\item[Insert the full typed Bellman ledger.]
Use Equations~\eqref{eq:timebellman}--\eqref{eq:spacebellman}.  Separate
one-time generation and acquisition in \(R_i\) from multiplicative
per-child-work overhead in \(w_e\).

\item[Insert the exact double-wrapper criterion.]
Use \eqref{eq:doubleaccept}, and use \eqref{eq:samescaleouter} for a
same-scale outer wrapper.  Mark the numerical \(17/10\) point
\textsc{Conditional}.

\item[Permit potential-collapsed administrative epochs.]
Add the mass theorem and the stack charge \(h+d\).  Require a locally
checkable potential, charged discovery, and a well-founded epoch rank.

\item[Add the same-realization guard.]
For input types, require Equation~\eqref{eq:samerealization}; do not use
\(\max_i(C_i+\rho_i)\) in place of \(\max_iC_i+\max_i\rho_i\).

\item[Add cycle hostile controls.]
Positive unit-product cycles have no finite budget.  Zero-weight unit-product
cycles may have fixed points but do not suppress the unknown seed exponent.
Type potentials cannot alter the restored cycle threshold.

\item[Preserve the open statement.]
The final missing theorem is an actual all-CNF mixed epoch or typed
double-wrapper realization with fully charged least budgets below
\(\gamma\).  None has yet been proved.
\end{description}

\section{Final status}

The universal source-only recursive cover was stronger than necessary.  The
correct relaxation is a universal, charged terminal-or-recursive epoch.  It
may be typed, may use direct \(\VALUE\) channels, and may amortize same-scale
administrative work through a source-checkable potential.  Once collapsed,
its recursive core must still contract on every feasible recurrent cycle,
and its global time and space maxima must be witnessed by one uniform
algorithm below the endpoint.

The full acceptance inequalities, including complete return cycles,
acquisition, replay, and live-space charges, are now explicit and
\textsc{Proved/Exact} as recurrence mathematics.  The two displayed favorable
points are \textsc{Conditional}.  The rational tests are
\textsc{Computational}.  The existence of a qualifying general-SAT compiler,
and therefore a defensible proof of \(P\ne NP\) from this route, remains
\textsc{Open}.

\begin{thebibliography}{9}

\bibitem{Williams2010}
R. Ryan Williams.
\newblock Alternation-trading proofs, linear programming, and lower bounds.
\newblock In \emph{27th International Symposium on Theoretical Aspects of
Computer Science}, LIPIcs 5, pages 669--680, 2010.
\newblock \url{https://doi.org/10.4230/LIPIcs.STACS.2010.2494}.

\bibitem{FortnowEtAl2005}
Lance Fortnow, Richard J. Lipton, Dieter van Melkebeek, and Anastasios Viglas.
\newblock Time-space lower bounds for satisfiability.
\newblock \emph{Journal of the ACM}, 52(6):835--865, 2005.
\newblock \url{https://doi.org/10.1145/1101821.1101822}.

\bibitem{FominGrandoniKratsch2009}
Fedor V. Fomin, Fabrizio Grandoni, and Dieter Kratsch.
\newblock A measure \& conquer approach for the analysis of exact algorithms.
\newblock \emph{Journal of the ACM}, 56(5), Article 25, 2009.
\newblock \url{https://doi.org/10.1145/1552285.1552286}.

\bibitem{MaehlmannSiebertzVigny2021}
Nikolas M\"ahlmann, Sebastian Siebertz, and Alexandre Vigny.
\newblock Recursive backdoors for SAT.
\newblock In \emph{46th International Symposium on Mathematical Foundations
of Computer Science}, LIPIcs 202, pages 73:1--73:18, 2021.
\newblock \url{https://doi.org/10.4230/LIPIcs.MFCS.2021.73}.

\bibitem{DreierOrdyniakSzeider2022}
Jan Dreier, Sebastian Ordyniak, and Stefan Szeider.
\newblock SAT backdoors: depth beats size.
\newblock In \emph{30th Annual European Symposium on Algorithms}, LIPIcs 244,
pages 46:1--46:18, 2022.
\newblock \url{https://doi.org/10.4230/LIPIcs.ESA.2022.46}.

\bibitem{OrdyniakSchidlerSzeider2021}
Sebastian Ordyniak, Andr\'e Schidler, and Stefan Szeider.
\newblock Backdoor DNFs.
\newblock In \emph{Proceedings of the Thirtieth International Joint
Conference on Artificial Intelligence}, pages 1403--1409, 2021.
\newblock \url{https://doi.org/10.24963/ijcai.2021/194}.

\bibitem{Williams2024}
R. Ryan Williams.
\newblock Self-improvement for circuit-analysis problems.
\newblock In \emph{Proceedings of the 56th Annual ACM Symposium on Theory of
Computing}, pages 1374--1385, 2024.
\newblock \url{https://doi.org/10.1145/3618260.3649723}.

\end{thebibliography}

% END CODEX REPORT::LEAD::FINAL-LINK2-ARC::2026-10-08
\end{document}
